\chapter{Opis założeń projektu}

\section{Cel projektu}
\label{sec:celProjektu}
Celem projektu jest stworzenie funkcjonalnego, przejrzystego i intuicyjnego systemu, który wspiera zarządzanie warsztatem samochodowym. Obejmuje to:

\begin{itemize}
    \item dodawanie danych klientów oraz ich pojazdów,
    \item tworzenie zleceń napraw z przypisanymi usługami, częściami i pracownikami,
    \item kontrolę i zarządzanie dostępnymi częściami i usługami,
    \item przeliczanie kosztów zlecenia na podstawie wybranych pozycji,
    \item zapewnienie prostego i przejrzystego interfejsu użytkownika.
\end{itemize}

\subsection{Analiza problemu}
\label{sec:analizaProblemu}
Zarządzanie warsztatem samochodowym to złożony proces, wymagający sprawnej organizacji informacji oraz płynnej obsługi klientów. Wiele warsztatów wciąż operuje na papierowej dokumentacji lub w niedostosowanych narzędziach ogólnego przeznaczenia, co wpływa negatywnie na efektywność pracy. Projekt ten odpowiada na realne potrzeby branży mechanicznej, oferując cyfrowe rozwiązanie dostosowane do pracy serwisu.
Problem ten wpływa bezpośrednio na efektywność pracy warsztatu oraz przejrzystość informacji. Istnieje zapotrzebowanie na prosty, ale funkcjonalny system, który umożliwiłby szybkie dodawanie zleceń, przypisywanie ich do konkretnych pracowników i klientów oraz kontrolowanie terminów i kosztów usług.

\subsection{Plan realizacji projektu}
\label{sec:planRealizacji}

Projekt zostanie wdrażany etapami, zgodnie z poniższą strukturą:
\begin{enumerate}
    \item Wstępna analiza potrzeb użytkowników oraz określenie kluczowych funkcjonalności.
    \item Stworzenie projektu relacyjnej bazy danych (MySQL), gwarantującej spójność i wydajność danych,
    \item Implementacja interfejsu graficznego z użyciem technologii JavaFX,
    \item Integracja logiki aplikacji z bazą danych,
    \item Przeprowadzenie testów w celu wyeliminowania błędów,
    \item Sporządzenie kompletnej dokumentacji obejmującej założenia projektu, opis architektury oraz instrukcję obsługi.
\end{enumerate}

Efektem finalnym będzie funkcjonalna aplikacja desktopowa wspomagająca działalność warsztatu samochodowego. Dodatkowo, system będzie posiadał zabezpieczenia oraz możliwość rozbudowy o dodatkowe funkcjonalności w przyszłości.


\section{Wymagania funkcjonalne}
\label{sec:wymaganiaFunkcjonalne}
System zarządzania serwisem samochodowym REPAIRO pozwala użytkownikom realizację konkretnych zadań w zależności od przypisanej roli. Każdy użytkownik, po zalogownaiu, ma dostęp wyłącznie do funkcjonalności niezbędnych do wykonywania swoich obowiązków. W tej sekcji opisano wymagania funkcjonalne procesu logowania a następnie przedstawiono podział funkcji dostępnych w aplikacji w zależności od roli użytkownika.

\subsection{Logowanie do systemu}
\label{sec:logowanie}

\begin{itemize}
    \item System umożliwia logowanie się do aplikacji i zapewnia dostęp wyłącznie uprawnionym użytkownikom. Uwierzytelnianie odbywa się na podstawie danych przechowywanych w bazie danych - użytkownik wprowadza nazwę użytkownika i hasło, a system sprawdza poprawność tych danych. W celu zapewnienia bezpieczeństwa, pole hasło domyślnie jest ukryte, z możliwością jego odsłonięcia za pomocą przycisku.
    \item System wyposażono w mechanizmy ochronne. W przypadku wykrycia nieznanej roli, użytkownik otrzymuje jednoznaczny komunikat o błędzie. Przy próbie logowania z błędnymi danymi (np. niestniejący użytkownik lub niepoprawne hasło), aplikacja wyświetla odpowiedni komunikat.
    \item W przypadku udanego logowania, system ustala rolę użytkownika (np. administrator, menadżer, mechanik, recepcjonista, magazynier) oraz pobiera dane z relacyjnych tabel bazy danych, takie jak pełna nazwa stanowiska oraz imię użytkownika, które są następnie wyświetlane na panelu głównym aplikacji. Interfejs dostowuje się do poziomu uprawnień każdej z ról, wyświetlając tylko dostępne widoki.
    \item Po zalogowaniu, dane użytkownika - jego imię, stanowisko oraz rola - są przechowywane w pamięci aplikacji aż do momentu jej zamknięcia.
\end{itemize}
Po udanym zalogowaniu, użytkownik zostaje przekierowany do panelu głównego, który jest dostosowany do jego roli. Każda rola ma przypisane odpowiednie funkcje związane z kompetencjami i obowiązkami użytkowników. Niezależnie od roli, każdy użytkownik ma dostęp do wspólnych widoków systemowych, takich jak ekran logowania, wylogowanie się czy widok informacji o aplikacji. Użytkownik nie widzi zakładek i widoków, do których nie ma przypisanych uprawnień. Poniżej opisano szczegółowo funkcje dostępne dla pięciu ról w systemie: administratora, menadżera, mechanika, recepcjonisty i magazyniera.

\subsection{Funkcjonalność dla roli administrator i menadżer}
\label{sec:funkcjonalnoscAdministratora}
Użytkownicy z rolą \textbf{administratora i menadżera} posiadają najwyższy stopień uprawnień w systemie. Logują się do systemu, wprowadzając unikalny login oraz hasło. Po zalogowaniu, mają dostęp do interfejsu, który umożliwia mu zarządzanie kluczowymi opcjami:

\begin{itemize}
    \item przeglądać, dodawać i usuwać pojazdy,
    \item przeglądać, dodawać i usuwać klientów,
    \item przeglądać, dodawać i usuwać części oraz usługi,
    \item przeglądać listę pracowników (na potrzeby projektu założono, że zarządzanie pracownikami odbywa się wyłącznie poprzez system bazodanowy),
    \item przeglądać, dodawać, edytować i usuwać zlecenia napraw,
    \item wyświetlać kalendarz wizyt, w którym prezentowane są aktywne zlecenia.
\end{itemize}

Administator ma również możliwość swobodnego przełączania się pomiędzy widokami aplikacji, co pozwala na łatwy dostęp do wszystkich funkcjonalności.

\subsection{Funkcjonalność dla roli mechanik}
\label{sec:funkcjonalnoscMechanika}
Użytkownicy z rolą \textbf{mechanika} logują się do systemu za pomocą unikalnego loginu i hasła. Po zalogowaniu, mają dostęp do interfejsu, który umożliwia mu zarządzanie kluczowymi opcjami.

Mechanik ma dostęp do następujących funkcji:
\begin{itemize}
    \item przeglądać, dodawać i usuwać pojazdy,
    \item przeglądać, dodawać i usuwać części oraz usługi,
    \item przeglądać, dodawać, edytować i usuwać zlecenia napraw,
    \item korzystania z dedykowanego widoku umożliwiającego sprawne przełączanie się między zakładkami i sekcjami aplikacji.
\end{itemize}

Interfejs graficzny został dostosowany do potrzeb pracownika technicznego - ograniczony do niezbędnych funkcjonalności, co upraszcza obsługę systemu.

\subsection{Funkcjonalność dla roli recepcjonista}
\label{sec:funkcjonalnoscRecepcjonisty}
Użytkownicy z rolą \textbf{recepcjonisty} logują się do systemu za pomocą unikalnego loginu i hasła. Po zalogowaniu, mają dostęp do interfejsu, który umożliwia mu zarządzanie kluczowymi opcjami.

Recepcjonista ma dostęp do następujących funkcji:
\begin{itemize}
    \item przeglądać, dodawać i usuwać pojazdy,
    \item przeglądać, dodawać i usuwać klientów,
    \item przeglądać, dodawać, edytować i usuwać zlecenia napraw,
    \item przeglądanie kalendarza wizyt w celu koordynacji pracy warsztatu.
\end{itemize}

Interfejs graficzny został dostosowany do potrzeb pracownika obsługi klienta - ograniczony do niezbędnych funkcjonalności, co upraszcza obsługę systemu.

\subsection{Funkcjonalność dla roli magazynier}
\label{sec:funkcjonalnoscMagazyniera}
Użytkownicy z rolą \textbf{magazyniera} logują się do systemu za pomocą unikalnego loginu i hasła. Po zalogowaniu, mają dostęp do interfejsu, który umożliwia mu zarządzanie kluczowymi opcjami.

Rola magazyniera ogranicza się jedynie do zarządzania częściami samochodowymi. W związku z tym, użytkownik ma dostęp jedynie do widoku części, gdzie może:
\begin{itemize}
    \item dodawać nowe części do systemu,
    \item przeglądać dostępne części,
    \item usuwać części, które nie są już potrzebne.
\end{itemize}

Interfejs graficzny został dostosowany do potrzeb pracownika magazynu - ograniczony do niezbędnych funkcjonalności, co upraszcza obsługę systemu.


Wszystkie role użytkowników mają również dostęp do wspólnych widoków systemowych, takich jak: ekran logowania, informacje o aplikacji, możliwość zmiany hasła oraz wylogowania się. Przełączanie między widokami jest dostępne zgodnie z przydzielonymi uprawnieniami - użytkownik widzi jedynie te opcje, do których posiada dostęp.

\subsection{Porównanie ról użytkowników i ich uprawnień}
\label{sec:podsumowanieFunkcjonalnosci}

Poniżej przedstawiono tabelę porównującą dostęp poszczególnych ról użytkowników do modułów systemu. Znakiem \checkmark oznaczono pełny dostęp, natomiast znakiem \ding{55} brak dostępu:

\begin{table}[H]
    \centering
    \caption{Dostęp poszczególnych ról użytkowników do modułów aplikacji.\label{tab:dostepy}}
    \begin{tabularx}{\linewidth}{l*{5}{>{\centering\arraybackslash}X}}
        \toprule
        \textbf{Moduł} & \textbf{Administrator} & \textbf{Manager} & \textbf{Mechanik} & \textbf{Recepcjonista} & \textbf{Magazynier} \\
        \midrule
        Samochody             & \checkmark & \checkmark & \checkmark & \checkmark & \ding{55} \\
        Klienci               & \checkmark & \checkmark & \ding{55}  & \checkmark & \ding{55} \\
        Przyjęcia             & \checkmark & \checkmark & \checkmark & \checkmark & \ding{55} \\
        Części                & \checkmark & \checkmark & \checkmark & \ding{55}  & \checkmark \\
        Wizyty (kalendarz)    & \checkmark & \checkmark & \ding{55}  & \checkmark & \ding{55} \\
        Pracownicy (podgląd)  & \checkmark & \checkmark & \ding{55}  & \ding{55}  & \ding{55} \\
        Usługi                & \checkmark & \checkmark & \checkmark & \ding{55}  & \ding{55} \\
        O aplikacji           & \checkmark & \checkmark & \checkmark & \checkmark & \checkmark \\
        \bottomrule
    \end{tabularx}
\end{table}

\section{Wymagania niefunkcjonalne}
\label{sec:wymaganiaNiefunkcjonalne}

\subsection{Środowisko testowe}
\label{sec:srodowiskoTestowe}
Testy aplikacji przeprowadzono w środowisku o konfiguracji:
\begin{itemize}
  \item Procesor: Intel Core i5-13500
  \item Karta graficzna: NVidia GeForce RTX 3060
  \item Pamięć RAM: 32 GB
  \item Dysk: SSD NVMe 1024 GB
  \item System operacyjny: Windows 11 Home
  \item IDE: IntelliJ IDEA 2024.3.3
  \item JDK: OpenJDK 23.0.2
  \item JavaFX: 23.0.1 (zgodnie z deklaracją xmlns)
  \item JDBC: mysql-connector-java-8.0.33
  \newpage
  \item Baza danych: MySQL (phpMyAdmin)
    \begin{itemize}
        \item Serwer: 127.0.0.1 via TCP/IP
        \item Typ serwera: MariaDB
        \item Połączenie z serwerem: SSL nie jest używany
        \item Wersja serwera: 10.4.32-MariaDB - mariadb.org binary distribution
        \item Wersja protokołu: 10
        \item Użytkownik: root@localhost
        \item Kodowanie znaków serwera: UTF-8 Unicode (utf8mb4)
    \end{itemize}
\end{itemize}

\subsection{Wydajność}
\label{sec:wydajnosc}
Aplikacja została zoptymalizowana pod kątem wydajności, aby zapewnić płynne działanie nawet przy dużej liczbie danych. Przeprowadzone testy wykazały:
\begin{itemize}
    \item Czas ładowania aplikacji wynosi średnio 2 sekundy.
    \item Interfejs użytkownika jest responsywny, a przełączanie między widokami odbywa się bez opóźnień.
    \item Operacje na bazie danych (dodawanie, edytowanie, usuwanie rekordów) są realizowane w czasie nieprzekraczającym 1 sekundy dla pojedynczych rekordów.
    \item W trakcie działania aplikacja zużywa średnio 350 MB pamięci RAM oraz nie obciąża procesora.
\end{itemize}

\subsection{Bezpieczeństwo}
\label{sec:bezpieczenstwo}

Aplikacja została wyposażona w mechanizmy zabezpieczeń, które mają na celu zapewnienie integralności i dostępności informacji.

\begin{itemize}
    \item Aplikacja za każdym razem wymaga logowania przed uzyskaniem dostępu do systemu.
    \item Każdy użytkownik posiada unikalną nazwę użytkownika,
    \item Pole hasła zostało ukryte, a użytkownik ma możliwość odsłonięcia go za pomocą przycisku.
    \item Uprawnienia użytkowników są kontrolowane na podstawie roli użytkownika,
    \item Formularze posiadają walidację danych - nie można zapisać pustych ani niepoprawnych wartości (np. litery w polu liczbowym),
    \item Komunikaty informujące o sukcesie lub błędzie wyświetlają jasną i klarowną użytkownikowi wiadomość,
    \item Pola wymagające powiązania z innymi rekordami (np. pojazd, klient) są wybierane z rozwijanych list, co eliminuje możliwość wpisania niepoprawnych danych.
\end{itemize}

Aplikacja została stworzona z myślą o dalszym rozwoju i może być łatwo rozszerzona o dodatkowe zabezpieczenia, takie jak szyfrowanie danych czy logowanie zdarzeń.

\subsection{Użyteczność}
\label{sec:uzytecznosc}
Użyteczność systemu odgrywa kluczową rolę w zapewnieniu efektywnej i przyjaznej obsługi przez użytkowników. Poniższe punkty definiują podstawowe założenia dotyczące wygody i intuicyjności korzystania z aplikacji.

\begin{itemize}
    \item Czytelny oraz intuicyjny interfejs z jasno opisynami przyciskami i polami edycji.
    \item Formularze z walidacją danych - użytkownik nie może wprowadzić błędnych danych, a system informuje o błędach w czasie rzeczywistym.
    \item Informacje zwrotne po wykonaniu operacji o sukcesie lub błędzie wyświetlane w formie komunikatów wyświetlane są natychmiast.
    \item Usuwanie rekordów poprzedzone jest potwierdzeniem, aby uniknąć przypadkowego usunięcia danych.
    \item Responsywność interfejsu - aplikacja działa płynnie, a przełączanie między widokami nie powoduje opóźnień.
    \item Użycie ComboBoxów i kalendarzy (DatePicker) do wyboru danych zamiast ręcznego wpisywania.
    \item Obsługa różnych ról użytkowników z dostosowanym widokiem (np. mechanik widzi inne funkcje niż administrator).
\end{itemize}

Przestrzeganie tych zasad ma na celu zwiększenie satysfakcji użytkowników oraz minimalizację błędów wynikających z nieporozumień lub trudności obsługi systemu.

\subsection{Skalowalność}
\label{sec:skalowalnosc}
Aplikacja została zaprojektowana z myślą o przyszłej rozbudowie i skalowalności. Poniższe aspekty zapewniają elastyczność i możliwość dostosowania systemu do rosnących potrzeb:
\begin{itemize}
    \item Możliwość dodania hashowania haseł w celu zapewnienia bezpieczeństwa danych użytkowników.
    \item Możliwość dodania dynamicznego zmieniania rozmiaru okna (powiększanie, pomniejszanie) poprawia użyteczność i pozwala na łatwiejszą adaptację interfejsu do różnych ekranów.
    \item Struktura bazy danych została zaprojektowana z myślą o rozbudowie, umożliwiająca dodawanie nowych tabel i relacji bez zaburzania działania istniejących funkcji.
    \item Interfejs aplikacji jest modularny i prosty do rozbudowy, co pozwala łatwo dodawać nowe zakładki i funkcjonalności bez przebudowy całego systemu.
    \item Możliwość rozbudowania aplikacji o integrację z innymi systemami, np. pocztą elektroniczną czy systemami płatności. 
    \item Struktura kodu i organizacja plików (jasny podział na warstwy, foldery według funkcjonalności i typów komponentów) ułatwiają orientację i rozwój aplikacji przez kolejnych programistów.
    \item Obecny model przypisywania części do zleceń jest ograniczony do jednego rodzaju części, jednak architektura pozwala na łatwe rozszerzenie tej funkcjonalności, by obsłużyć wiele części dla jednego zlecenia.
\end{itemize}

Zorganizowana struktura kodu oraz bazy danych umożliwia łatwe wprowadzanie zmian i dodawanie nowych funkcji, co jest kluczowe w kontekście dynamicznie zmieniających się potrzeb użytkowników i rynku.

\subsection{Integralność danych}
\label{sec:integralnoscDanych}

Integralność danych gwarantuje, że informacje przechowywane i przetwarzane w systemie są spójne, poprawne i wiarygodne przez cały czas działania aplikacji. 

\begin{itemize}
    \item \textbf{Klucze główne i obce:} - odpowiednio zdefiniowane relacje między tabelami w bazie danych zapobiegają powstawaniu niespójnych rekordów.
    \item \textbf{Walidacja danych na poziomie aplikacji} - aplikacja sprawdza poprawność wprowadzanych danych przed ich zapisaniem do bazy, co minimalizuje ryzyko błędów.
    \item \textbf{Spójność danych wprowadzanych przez użytkownika} - system wymaga uzupełnienia wszystkich niezbędnych informacji, eliminując możliwość pozostawienia pustych lub błędnych pól.
    \item \textbf{Obsługa błędów} - system odpowiednio reaguje na błędy podczas operacji na danych, wyświetlając użytkownikowi klarowne komunikaty o sukcesie lub porażce.
\end{itemize}

Zastosowanie tych rozwiązań gwarantuje spójność danych w systemie oraz minimalizuje ryzyko błędów, które mogłyby negatywnie wpłynąć na funkcjonowanie aplikacji.